
\section{Abel's Method of Rationalizing a Ratio} As an application of Sylvester's Resultant, let us apply the theorem to the case of two polynomials which are related to Abel's Proof of the 
insolubility of the quintic equation in terms of radicals. In particular let us examine Abel's method of rationalizing a polynomial function of radicals.
Let 
\begin{align*}
f(x)  &= x^{p} - R\\
g(x) &=  a_{p-1} x^{p-1}  + a_{p-2} x^{p-2} + \hdots + a_2 x^2 + a_1 x   + a_0,  \\
\end{align*} where $R$ is not a $p$th power of an element in the field of coefficients. 
$f(x)$ and $g(x)$ will have a common root if the Sylvester Resultant formed by these two functions vanishes. 
The roots of $f(x) = 0$ take the form $z_k = \omega^k R^{1\over p}$ and we're assuming that $R^{1\over p}$ is an irrationality in the domain containing the coefficients. 
Also with respect to Abel's rationalization method 
we need to show that 
\be \prod_{k=0}^{p-1}\, g(z_k) = G(z^p)  = G(R)  \mbox{  where $G(R)$ is rational in the domain of the coefficients of } g(x).\ee
Using $f(x)$ and $g(x)$ we construct the Sylvester Resultant
\[\det(M) = \prod_{k=0}^{p-1}\, g(z_k)  =
\det \begin{bmatrix}
1     & 0 & \cdots  & -R                   &            &    \\
 0    &  1 & \ddots & \cdots               & \ddots  &    \\
 \vdots  &  & \ddots & \cdots               & \ddots  &    \\
 0      &  \hdots   &            & 1                       & \cdots  & -R\\
a_{p-1} &     & \cdots & ~~~~~~a_0 &             & \\
       &     & \ddots & \cdots                & \ddots & \\
       &     &            &a_{p-1}                     & \cdots &a_0 \\
\end{bmatrix}.\]

Since the roots of $f(x) = x^p - R $ are exactly $z_k = \omega^k R^{1\over p}$, substituting them into the formula yields 
\[ \det(M) =  \prod_{k=0}^{p-1}\, g(z_k) .\] 
Because computing the determinant $\det(M)$ only involves additions, subtractions, and multiplication of the matrix entries, the 
final polynomial expression outputted by the determinant can only contain powers of combinations of the terms in the matrix M. 
Therefore the resulting polynomial must be a function of $x^p = R$ and the coefficients and must therefore be a rational expression. 
The above is a demonstration of Abel's method of Rationalization. 